% Options for packages loaded elsewhere
% Options for packages loaded elsewhere
\PassOptionsToPackage{unicode}{hyperref}
\PassOptionsToPackage{hyphens}{url}
\PassOptionsToPackage{dvipsnames,svgnames,x11names}{xcolor}
%
\documentclass[
  letterpaper,
  DIV=11,
  numbers=noendperiod]{scrartcl}
\usepackage{xcolor}
\usepackage{amsmath,amssymb}
\setcounter{secnumdepth}{-\maxdimen} % remove section numbering
\usepackage{iftex}
\ifPDFTeX
  \usepackage[T1]{fontenc}
  \usepackage[utf8]{inputenc}
  \usepackage{textcomp} % provide euro and other symbols
\else % if luatex or xetex
  \usepackage{unicode-math} % this also loads fontspec
  \defaultfontfeatures{Scale=MatchLowercase}
  \defaultfontfeatures[\rmfamily]{Ligatures=TeX,Scale=1}
\fi
\usepackage{lmodern}
\ifPDFTeX\else
  % xetex/luatex font selection
\fi
% Use upquote if available, for straight quotes in verbatim environments
\IfFileExists{upquote.sty}{\usepackage{upquote}}{}
\IfFileExists{microtype.sty}{% use microtype if available
  \usepackage[]{microtype}
  \UseMicrotypeSet[protrusion]{basicmath} % disable protrusion for tt fonts
}{}
\makeatletter
\@ifundefined{KOMAClassName}{% if non-KOMA class
  \IfFileExists{parskip.sty}{%
    \usepackage{parskip}
  }{% else
    \setlength{\parindent}{0pt}
    \setlength{\parskip}{6pt plus 2pt minus 1pt}}
}{% if KOMA class
  \KOMAoptions{parskip=half}}
\makeatother
% Make \paragraph and \subparagraph free-standing
\makeatletter
\ifx\paragraph\undefined\else
  \let\oldparagraph\paragraph
  \renewcommand{\paragraph}{
    \@ifstar
      \xxxParagraphStar
      \xxxParagraphNoStar
  }
  \newcommand{\xxxParagraphStar}[1]{\oldparagraph*{#1}\mbox{}}
  \newcommand{\xxxParagraphNoStar}[1]{\oldparagraph{#1}\mbox{}}
\fi
\ifx\subparagraph\undefined\else
  \let\oldsubparagraph\subparagraph
  \renewcommand{\subparagraph}{
    \@ifstar
      \xxxSubParagraphStar
      \xxxSubParagraphNoStar
  }
  \newcommand{\xxxSubParagraphStar}[1]{\oldsubparagraph*{#1}\mbox{}}
  \newcommand{\xxxSubParagraphNoStar}[1]{\oldsubparagraph{#1}\mbox{}}
\fi
\makeatother


\usepackage{longtable,booktabs,array}
\newcounter{none} % for unnumbered tables
\usepackage{calc} % for calculating minipage widths
% Correct order of tables after \paragraph or \subparagraph
\usepackage{etoolbox}
\makeatletter
\patchcmd\longtable{\par}{\if@noskipsec\mbox{}\fi\par}{}{}
\makeatother
% Allow footnotes in longtable head/foot
\IfFileExists{footnotehyper.sty}{\usepackage{footnotehyper}}{\usepackage{footnote}}
\makesavenoteenv{longtable}
\usepackage{graphicx}
\makeatletter
\newsavebox\pandoc@box
\newcommand*\pandocbounded[1]{% scales image to fit in text height/width
  \sbox\pandoc@box{#1}%
  \Gscale@div\@tempa{\textheight}{\dimexpr\ht\pandoc@box+\dp\pandoc@box\relax}%
  \Gscale@div\@tempb{\linewidth}{\wd\pandoc@box}%
  \ifdim\@tempb\p@<\@tempa\p@\let\@tempa\@tempb\fi% select the smaller of both
  \ifdim\@tempa\p@<\p@\scalebox{\@tempa}{\usebox\pandoc@box}%
  \else\usebox{\pandoc@box}%
  \fi%
}
% Set default figure placement to htbp
\def\fps@figure{htbp}
\makeatother





\setlength{\emergencystretch}{3em} % prevent overfull lines

\providecommand{\tightlist}{%
  \setlength{\itemsep}{0pt}\setlength{\parskip}{0pt}}



 


\KOMAoption{captions}{tableheading}
\makeatletter
\@ifpackageloaded{caption}{}{\usepackage{caption}}
\AtBeginDocument{%
\ifdefined\contentsname
  \renewcommand*\contentsname{Table of contents}
\else
  \newcommand\contentsname{Table of contents}
\fi
\ifdefined\listfigurename
  \renewcommand*\listfigurename{List of Figures}
\else
  \newcommand\listfigurename{List of Figures}
\fi
\ifdefined\listtablename
  \renewcommand*\listtablename{List of Tables}
\else
  \newcommand\listtablename{List of Tables}
\fi
\ifdefined\figurename
  \renewcommand*\figurename{Figure}
\else
  \newcommand\figurename{Figure}
\fi
\ifdefined\tablename
  \renewcommand*\tablename{Table}
\else
  \newcommand\tablename{Table}
\fi
}
\@ifpackageloaded{float}{}{\usepackage{float}}
\floatstyle{ruled}
\@ifundefined{c@chapter}{\newfloat{codelisting}{h}{lop}}{\newfloat{codelisting}{h}{lop}[chapter]}
\floatname{codelisting}{Listing}
\newcommand*\listoflistings{\listof{codelisting}{List of Listings}}
\makeatother
\makeatletter
\makeatother
\makeatletter
\@ifpackageloaded{caption}{}{\usepackage{caption}}
\@ifpackageloaded{subcaption}{}{\usepackage{subcaption}}
\makeatother
\usepackage{bookmark}
\IfFileExists{xurl.sty}{\usepackage{xurl}}{} % add URL line breaks if available
\urlstyle{same}
% fallback for those not using the hyperref driver hyperxmp:
\makeatletter
\@ifundefined{xmpquote}{\newcommand{\xmpquote}[1]{#1}}{}
\makeatother
\hypersetup{
  pdftitle={Reproducibility Report for Study by De la Fuente et al.~(2021{,} Psychological Science)},
  pdfauthor={Fei Xu},
  colorlinks=true,
  linkcolor={blue},
  filecolor={Maroon},
  citecolor={Blue},
  urlcolor={Blue},
  pdfcreator={LaTeX via pandoc}}


\title{Reproducibility Report for Study by De la Fuente et al.~(2021,
Psychological Science)}
\author{Fei Xu}
\date{2026-10-03}
\begin{document}
\maketitle

\renewcommand*\contentsname{Table of contents}
{
\setcounter{tocdepth}{3}
\tableofcontents
}

\subsection{Introduction}\label{introduction}

{[}No abstract is needed.{]} Each reproducibility project will have a
straightforward, no frills report of the study and reproducibility
results. These reports will be publicly available as supplementary
material for the aggregate report(s) of the project as a whole. Also, to
maximize project integrity, the intro and methods will be written and
critiqued in advance of your attempt to reproduce the results.
Introductions can be just 1-2 paragraphs clarifying the main idea of the
original study, the target finding for your reproducibility attempt, and
any other essential information. It will NOT have a literature review --
that is in the original publication. You can write both the introduction
and the methods in past tense.

\textbf{Clarify key analysis of interest here} You can also pre-specify
additional analyses you plan to do.

\subsubsection{Justification for choice of
study}\label{justification-for-choice-of-study}

De la Fuente et al.~(2021) investigates how wild marmosets adjust their
foraging partner choices and strength of foraging associations with
their conspecifics under varied food availability. The authors
manipulated the abundance and distribution of food resources to look at
how the marmosets of different sex, hierarchy, and age interact with
each other. Among many of their hypothesis, I would like to focus on two
of them that involve juveniles: (a)social tolerance towards juveniles is
high as juveniles would share a feeding site with a greater number of
group members and for greater proportion of time than adults, and (b)
dyads composed of an adult and a juvenile will spend greater proportions
of time sharing a feeding site than dyads composed of two adults. These
two hypothesis align with my broad research interest on marmoset
development, focusing particularly how different
parent-infant/helper-offspring interactions contribute to offspring
behavioral and brain development. The adult-juvenile interaction from
this study can provide me with insights into the flexibility of social
tolerance in marmosets. In addition, the study manipulates a large
dataset, and simulating the data allows me to learn how to deal with
huge database with a wide range of predictor variables.

\subsubsection{Anticipated challenges}\label{anticipated-challenges}

The original study contains 5 food availability conditions and 2
dataset(individual and dyads). Even though I chose to focus on only two
of the many hypothesis, simulation would require me to simulate both
datasets. Taking the individual dataset as an example, I will first
calculate the node strength and node degree from the original dataset
and refit the original GLMMs with the models presented in the paper.
Next, I will keep all the original variables (Condition, Session, Group,
Individual, Sex, Age, Rank, Food\_consumed, Avoidance\_events, and
Agression\_events) for each observation and use the refitted models to
predict the expected response for each observation. I will generate node
strength values from the predicted binomial probability and node degree
values from the predicted Poisson means. Because the authors didn't
provide the original code, the reconstructed models may not reproduce
the exact original implementation.

\subsubsection{Links}\label{links}

Project repository (on Github):
https://github.com/FeiXu27/DelaFuente2021

Original paper (as hosted in your repo):
https://github.com/FeiXu27/DelaFuente2021/blob/main/original\_paper/Journal\%20of\%20Animal\%20Ecology\%20-\%202021\%20-\%20De\%20la\%20Fuente\%20-\%20Foraging\%20networks\%20and\%20social\%20tolerance\%20in\%20a\%20cooperatively\%20breeding.pdf

\subsection{Methods}\label{methods}

\subsubsection{Description of the steps required to reproduce the
results}\label{description-of-the-steps-required-to-reproduce-the-results}

Please describe all the steps necessary to reproduce the key result(s)
of this study.

\subsubsection{Differences from original
study}\label{differences-from-original-study}

Explicitly describe known differences in the analysis pipeline between
the original paper and yours (e.g., computing environment). The goal, of
course, is to minimize those differences, but differences may occur.
Also, note whether such differences are anticipated to influence your
ability to reproduce the original results.

\subsection{Project Progress Check 1}\label{project-progress-check-1}

\subsubsection{Measure of success}\label{measure-of-success}

Please describe the outcome measure for the success or failure of your
reproduction and how this outcome will be computed.

\subsubsection{Pipeline progress}\label{pipeline-progress}

Earlier in this report, you described the steps necessary to reproduce
the key result(s) of this study. Please describe your progress on each
of these steps (e.g., data preprocessing, model fitting, model
evaluation).

\subsection{Results}\label{results}

\subsubsection{Data preparation}\label{data-preparation}

Data preparation following the analysis plan.

\subsubsection{Key analysis}\label{key-analysis}

The analyses as specified in the analysis plan.

\emph{Side-by-side graph with original graph is ideal here}

\subsubsection{Exploratory analyses}\label{exploratory-analyses}

Any follow-up analyses desired (not required).

\subsection{Discussion}\label{discussion}

\subsubsection{Summary of Reproduction
Attempt}\label{summary-of-reproduction-attempt}

Open the discussion section with a paragraph summarizing the primary
result from the key analysis and assess whether you successfully
reproduced it, partially reproduced it, or failed to reproduce it.

\subsubsection{Commentary}\label{commentary}

Add open-ended commentary (if any) reflecting (a) insights from
follow-up exploratory analysis of the dataset, (b) assessment of the
meaning of the successful or unsuccessful reproducibility attempt -
e.g., for a failure to reproduce the original findings, are the
differences between original and present analyses ones that definitely,
plausibly, or are unlikely to have been moderators of the result, and
(c) discussion of any objections or challenges raised by the current and
original authors about the reproducibility attempt (if you contacted
them). None of these need to be long.




\end{document}
